\section{Instrument: harness, restoration surface, and benchmark tiers}
\label{sec:evaluation}

\emph{Status.} This appendix specifies the instrument that \emph{would} realise
the causal contract; the benchmark tiers, fault families and gates below are a
retained registered plan and none of them has been executed on a live agent.  The
policy-ranking study named in the tiers is deferred, and the instrument gate is
unopened.  The causal contract of \cref{sec:contract,sec:causal} is
benchmark-independent, and each requirement below is written so that it is
checkable before any scientific run.

\subsection{Agent harness}

The signed R5-2A record binds the controlled tabular runtime, model, and
environment.  Pi remains non-qualifying until a replacement audit and signature
bind its exact stack; current Pi runs cannot count as E0.  The deterministic
fixture is likewise not scientific.  In the
harness-optimization literature of \cref{sec:related} the harness is the object
being edited; here it is pinned precisely because it is \emph{not} the object of
study --- the intervention is the one-step regime at a checkpoint, not the
runtime that carries it.  A second local model family may be declared as a
generalisation probe and would be reported separately, never pooled.  Workload
seeds, sampler settings, and model randomness are recorded, and deterministic
behavior is claimed only if the production audit demonstrates it.

\subsection{The restoration surface}

Branching requires that every component of the checkpoint state be restored, not
merely the parts that are convenient to snapshot.  \Cref{tab:restoration}
enumerates the surface and, for each layer, the mechanism and its current
manuscript-admitted evidential status: \textbf{D} documented by an authoritative source,
\textbf{T} demonstrated only in this project's deterministic scaffold,
\textbf{U} unverified by a qualifying E0.  Local smoke, profiling, or a
non-qualifying capability audit cannot promote these labels.

\begin{table}[ht]
\centering\small
\begin{tabular}{@{}p{3.3cm}p{6.3cm}p{2.1cm}c@{}}
\toprule
Layer & Mechanism & Failure mode & Status \\
\midrule
Workspace filesystem & Content-addressed snapshot and restore of the task tree & Silent partial restore & T \\
Executable environment & Container image digest pinned; dependency lock recorded & Drift across branches & U \\
Agent session state & Controlled-agent history and sampler configuration replayed into the pinned local runtime & Truncated history or configuration drift & U \\
Budget counters & Token, spend, and wall-clock ledger restored to the checkpoint value & Branch inherits sibling spend & T \\
External tool state & Declared inside or outside the experimental boundary; anything inside must be snapshottable & Undeclared shared state & U \\
Local inference cache & Disabled or isolated by branch in the production adapter; never treated as restored state & Cross-branch dependence or branch-position effects & U \\
Sampling randomness & Local seed and sampler settings recorded; deterministic support tested rather than assumed & False determinism claim & U \\
\bottomrule
\end{tabular}
\caption{Restoration surface.  A checkpoint is eligible only when every layer is
demonstrated; local cache state is disabled or isolated and tested by branch
position rather than claimed to be restored (\cref{sec:cache}).}
\label{tab:restoration}
\end{table}

\subsection{Cache accounting for local and API-backed execution}
\label{sec:cache}

Sibling branches share the trajectory prefix by construction.  In an API-backed
deployment with a provider cache, branch order can therefore change billed cost
for identical work.  Under the illustrative registered price vector, the
first-versus-cache-warm ratio is $3.6\times$, $5.9\times$, and $8.2\times$ over
the three $50$k--$600$k prefix scenarios.  These are design calculations for
API-backed deployments, not an operating range of the registered study, whose
production candidate uses pinned local inference and local metering.

This has two consequences the protocol must absorb.

\emph{Cost.}  The registered constraint \eqref{eq:cost-constraint} is evaluated
on cache-neutralised resource components and the local meter records the
corresponding runtime use.  Provider-billed reconciliation is not applicable to
this substrate.  A later hosted-endpoint adapter would require a new approval,
a billing-precision tolerance, and provider-cache reconciliation before E0.

\emph{Outcomes.}  The production adapter disables or branch-isolates local
inference caches and E0 still compares outcomes across branch positions.  A
detected cache or position dependence is an interference finding that invalidates
the isolation premise of \cref{sec:causal}, not a nuisance to adjust away.

\subsection{Benchmark tiers}

The tiers carry distinct, non-interchangeable roles, reported per tier rather
than pooled; Tier~4 is a declared negative control on executability
(\cref{tab:tiers}, \cref{app:protocol}).  Tier~1 has two components: the
primary F2-FEHPO substrate, confirmatory only under its active R5-2A signature;
the current controlled-runtime signature excludes Pi.  The earlier MLGym F1-ORD and F2-ADV calculations are historical
substrate-specific contrasts only; they are not executed in the current protocol
and cannot establish, gate, replicate, or rescue F2-FEHPO.  Tier~0 is a
non-qualifying infrastructure smoke test and never validates E0.

Licensing is a design constraint rather than a footnote.  MLGym is excluded from
the current protocol; any future use or redistribution requires a new scope and
licence decision, and task-level dataset terms are checked separately.  Public tasks additionally
create contamination risk, requiring version pinning, disclosure of model
training-date limitations, and avoidance of protected evaluation material.
MLE-bench Lite is held out until after the confirmatory freeze precisely so that
its grading infrastructure remains uncontaminated by design iteration.

\subsection{Fault families}

Three roles are separated because \cref{thm:ordinal} requires the primary
positive population to arise independently of the risk score being tested.

\textbf{F1-ORD, historical provenance/dependency contrast.}  The earlier MLGym
calculation substituted a dependency, manifest, or artefact with a
machine-checkable postcondition.  It is not ported to the controlled substrate
and no formal equivalence tolerance is active; it remains a non-transferable
historical design diagnostic.  Because the authors saw the F1 precursor result
before proposing $\delta_{F1}=0.01$, no blindness or outcome-blindness claim is
made; \cref{app:faults} records the retired control.

\textbf{F2-FEHPO, primary randomized positive condition.}  An eligible feature
pipeline appears weak under a mismatched or inherited model configuration but
becomes valuable under a registered matched configuration.  Eligibility and
randomized matched-versus-mismatched assignment use pre-treatment state and
never read the risk score.  CoFEH documents the motivating greedy-myopia
mechanism in joint feature engineering and hyperparameter optimization
\citep{cofeh2026}; whether it occurs at these agent checkpoints remains an
empirical question.

\textbf{F2-ADV, historical adversarial-stress contrast.}  The previous
lower-risk-tercile MLGym calculations describe an intentionally anti-ordered
population.  They are not ported or executed in the current protocol and cannot
establish or rescue the primary ordinal-insufficiency claim.

\Cref{app:faults} gives the frozen machine-readable semantics and records the
historical simulation without treating it as the operative power calculation.
The F2-FEHPO pair, E3 root count, and repetition rule are hash-bound; R5-2B
selects the count mechanically after a disjoint pilot.  Runtime replacement
requires re-binding before E0.
